\documentclass[
  doc,
  floatsintext,
  longtable,
  a4paper,
  nolmodern,
  notxfonts,
  notimes,
  12pt,
  colorlinks=true,linkcolor=blue,citecolor=blue,urlcolor=blue]{apa7}

\usepackage{amsmath}
\usepackage{amssymb}

\geometry{inner=1in, outer=1in}
\fancyhfoffset[LE,RO]{0cm}


\usepackage[bidi=default]{babel}
\babelprovide[main,import]{spanish}


% get rid of language-specific shorthands (see #6817):
\let\LanguageShortHands\languageshorthands
\def\languageshorthands#1{}

\RequirePackage{longtable}
\RequirePackage{threeparttablex}

\makeatletter
\renewcommand{\paragraph}{\@startsection{paragraph}{4}{\parindent}%
	{0\baselineskip \@plus 0.2ex \@minus 0.2ex}%
	{-.5em}%
	{\normalfont\normalsize\bfseries\typesectitle}}

\renewcommand{\subparagraph}[1]{\@startsection{subparagraph}{5}{0.5em}%
	{0\baselineskip \@plus 0.2ex \@minus 0.2ex}%
	{-\z@\relax}%
	{\normalfont\normalsize\bfseries\itshape\hspace{\parindent}{#1}\textit{\addperi}}{\relax}}
\makeatother




\usepackage{longtable, booktabs, multirow, multicol, colortbl, hhline, caption, array, float, xpatch}
\usepackage{subcaption}


\renewcommand\thesubfigure{\Alph{subfigure}}
\setcounter{topnumber}{2}
\setcounter{bottomnumber}{2}
\setcounter{totalnumber}{4}
\renewcommand{\topfraction}{0.85}
\renewcommand{\bottomfraction}{0.85}
\renewcommand{\textfraction}{0.15}
\renewcommand{\floatpagefraction}{0.7}

\usepackage{tcolorbox}
\tcbuselibrary{listings,theorems, breakable, skins}
\usepackage{fontawesome5}

\definecolor{quarto-callout-color}{HTML}{909090}
\definecolor{quarto-callout-note-color}{HTML}{0758E5}
\definecolor{quarto-callout-important-color}{HTML}{CC1914}
\definecolor{quarto-callout-warning-color}{HTML}{EB9113}
\definecolor{quarto-callout-tip-color}{HTML}{00A047}
\definecolor{quarto-callout-caution-color}{HTML}{FC5300}
\definecolor{quarto-callout-color-frame}{HTML}{ACACAC}
\definecolor{quarto-callout-note-color-frame}{HTML}{4582EC}
\definecolor{quarto-callout-important-color-frame}{HTML}{D9534F}
\definecolor{quarto-callout-warning-color-frame}{HTML}{F0AD4E}
\definecolor{quarto-callout-tip-color-frame}{HTML}{02B875}
\definecolor{quarto-callout-caution-color-frame}{HTML}{FD7E14}

%\newlength\Oldarrayrulewidth
%\newlength\Oldtabcolsep


\usepackage{hyperref}



\usepackage{color}
\usepackage{fancyvrb}
\newcommand{\VerbBar}{|}
\newcommand{\VERB}{\Verb[commandchars=\\\{\}]}
\DefineVerbatimEnvironment{Highlighting}{Verbatim}{commandchars=\\\{\}}
% Add ',fontsize=\small' for more characters per line
\usepackage{framed}
\definecolor{shadecolor}{RGB}{241,243,245}
\newenvironment{Shaded}{\begin{snugshade}}{\end{snugshade}}
\newcommand{\AlertTok}[1]{\textcolor[rgb]{0.68,0.00,0.00}{#1}}
\newcommand{\AnnotationTok}[1]{\textcolor[rgb]{0.37,0.37,0.37}{#1}}
\newcommand{\AttributeTok}[1]{\textcolor[rgb]{0.40,0.45,0.13}{#1}}
\newcommand{\BaseNTok}[1]{\textcolor[rgb]{0.68,0.00,0.00}{#1}}
\newcommand{\BuiltInTok}[1]{\textcolor[rgb]{0.00,0.23,0.31}{#1}}
\newcommand{\CharTok}[1]{\textcolor[rgb]{0.13,0.47,0.30}{#1}}
\newcommand{\CommentTok}[1]{\textcolor[rgb]{0.37,0.37,0.37}{#1}}
\newcommand{\CommentVarTok}[1]{\textcolor[rgb]{0.37,0.37,0.37}{\textit{#1}}}
\newcommand{\ConstantTok}[1]{\textcolor[rgb]{0.56,0.35,0.01}{#1}}
\newcommand{\ControlFlowTok}[1]{\textcolor[rgb]{0.00,0.23,0.31}{\textbf{#1}}}
\newcommand{\DataTypeTok}[1]{\textcolor[rgb]{0.68,0.00,0.00}{#1}}
\newcommand{\DecValTok}[1]{\textcolor[rgb]{0.68,0.00,0.00}{#1}}
\newcommand{\DocumentationTok}[1]{\textcolor[rgb]{0.37,0.37,0.37}{\textit{#1}}}
\newcommand{\ErrorTok}[1]{\textcolor[rgb]{0.68,0.00,0.00}{#1}}
\newcommand{\ExtensionTok}[1]{\textcolor[rgb]{0.00,0.23,0.31}{#1}}
\newcommand{\FloatTok}[1]{\textcolor[rgb]{0.68,0.00,0.00}{#1}}
\newcommand{\FunctionTok}[1]{\textcolor[rgb]{0.28,0.35,0.67}{#1}}
\newcommand{\ImportTok}[1]{\textcolor[rgb]{0.00,0.46,0.62}{#1}}
\newcommand{\InformationTok}[1]{\textcolor[rgb]{0.37,0.37,0.37}{#1}}
\newcommand{\KeywordTok}[1]{\textcolor[rgb]{0.00,0.23,0.31}{\textbf{#1}}}
\newcommand{\NormalTok}[1]{\textcolor[rgb]{0.00,0.23,0.31}{#1}}
\newcommand{\OperatorTok}[1]{\textcolor[rgb]{0.37,0.37,0.37}{#1}}
\newcommand{\OtherTok}[1]{\textcolor[rgb]{0.00,0.23,0.31}{#1}}
\newcommand{\PreprocessorTok}[1]{\textcolor[rgb]{0.68,0.00,0.00}{#1}}
\newcommand{\RegionMarkerTok}[1]{\textcolor[rgb]{0.00,0.23,0.31}{#1}}
\newcommand{\SpecialCharTok}[1]{\textcolor[rgb]{0.37,0.37,0.37}{#1}}
\newcommand{\SpecialStringTok}[1]{\textcolor[rgb]{0.13,0.47,0.30}{#1}}
\newcommand{\StringTok}[1]{\textcolor[rgb]{0.13,0.47,0.30}{#1}}
\newcommand{\VariableTok}[1]{\textcolor[rgb]{0.07,0.07,0.07}{#1}}
\newcommand{\VerbatimStringTok}[1]{\textcolor[rgb]{0.13,0.47,0.30}{#1}}
\newcommand{\WarningTok}[1]{\textcolor[rgb]{0.37,0.37,0.37}{\textit{#1}}}

\providecommand{\tightlist}{%
  \setlength{\itemsep}{0pt}\setlength{\parskip}{0pt}}
\usepackage{longtable,booktabs,array}
\newcounter{none} % for unnumbered tables
\usepackage{calc} % for calculating minipage widths
% Correct order of tables after \paragraph or \subparagraph
\usepackage{etoolbox}
\makeatletter
\patchcmd\longtable{\par}{\if@noskipsec\mbox{}\fi\par}{}{}
\makeatother
% Allow footnotes in longtable head/foot
\IfFileExists{footnotehyper.sty}{\usepackage{footnotehyper}}{\usepackage{footnote}}
\makesavenoteenv{longtable}

\usepackage{graphicx}
\makeatletter
\newsavebox\pandoc@box
\newcommand*\pandocbounded[1]{% scales image to fit in text height/width
  \sbox\pandoc@box{#1}%
  \Gscale@div\@tempa{\textheight}{\dimexpr\ht\pandoc@box+\dp\pandoc@box\relax}%
  \Gscale@div\@tempb{\linewidth}{\wd\pandoc@box}%
  \ifdim\@tempb\p@<\@tempa\p@\let\@tempa\@tempb\fi% select the smaller of both
  \ifdim\@tempa\p@<\p@\scalebox{\@tempa}{\usebox\pandoc@box}%
  \else\usebox{\pandoc@box}%
  \fi%
}
% Set default figure placement to htbp
\def\fps@figure{htbp}
\makeatother







\usepackage{newtx}

\defaultfontfeatures{Scale=MatchLowercase}
\defaultfontfeatures[\rmfamily]{Ligatures=TeX,Scale=1}





\title{Introducción a la programación con Python: Guía rigurosa basada
en la documentación oficial de Python --- sin variables, tipos,
condicionales ni funciones}


\shorttitle{INTRODUCCION A PYTHON}


\usepackage{etoolbox}



\ccoppy{\textcopyright~2021}





\authorsnames{Edison Achalma,Edison Achalma}





\affiliation{
{Escuela Profesional de Economía, Universidad Nacional de San Cristóbal
de Huamanga}}




\leftheader{Achalma and Achalma}

\date{2026-06-21}


\abstract{Este abstract será actualizado una vez que se complete el
contenido final del artículo. }

\keywords{keyword1, keyword2}

\authornote{\par{\addORCIDlink{Edison Achalma}{0000-0001-6996-3364}} 
\par{ }
\par{   El autor no tiene conflictos de interés que revelar.    Los
roles de autor se clasificaron utilizando la taxonomía de roles de
colaborador (CRediT; https://credit.niso.org/) de la siguiente
manera:  Edison Achalma:   conceptualización, metodología, análisis
formal, investigación, recursos, curación de
datos, redacción, visualización, supervisión, administración del
proyecto}
\par{La correspondencia relativa a este artículo debe dirigirse a Edison
Achalma, Escuela Profesional de Economía, Universidad Nacional de San
Cristóbal de Huamanga, Portal Independencia N°
57, Ayacucho, AYA 5001, Perú, Email: \href{mailto:elmer.achalma.09@unsch.edu.pe}{elmer.achalma.09@unsch.edu.pe}}
}

\makeatletter
\let\endoldlt\endlongtable
\def\endlongtable{
\hline
\endoldlt
}
\makeatother

\urlstyle{same}



\hypersetup{
  pdftitle={Introducci\'{o}n a la programaci\'{o}n con Python},
  pdfauthor={Autor}
}
\makeatletter
\@ifpackageloaded{caption}{}{\usepackage{caption}}
\AtBeginDocument{%
\ifdefined\contentsname
  \renewcommand*\contentsname{Tabla de contenidos}
\else
  \newcommand\contentsname{Tabla de contenidos}
\fi
\ifdefined\listfigurename
  \renewcommand*\listfigurename{Lista de Figuras}
\else
  \newcommand\listfigurename{Lista de Figuras}
\fi
\ifdefined\listtablename
  \renewcommand*\listtablename{Lista de Tablas}
\else
  \newcommand\listtablename{Lista de Tablas}
\fi
\ifdefined\figurename
  \renewcommand*\figurename{Figura}
\else
  \newcommand\figurename{Figura}
\fi
\ifdefined\tablename
  \renewcommand*\tablename{Tabla}
\else
  \newcommand\tablename{Tabla}
\fi
}
\@ifpackageloaded{float}{}{\usepackage{float}}
\floatstyle{ruled}
\@ifundefined{c@chapter}{\newfloat{codelisting}{h}{lop}}{\newfloat{codelisting}{h}{lop}[chapter]}
\floatname{codelisting}{Listado}
\newcommand*\listoflistings{\listof{codelisting}{Lista de Listados}}
\makeatother
\makeatletter
\makeatother
\makeatletter
\@ifpackageloaded{caption}{}{\usepackage{caption}}
\@ifpackageloaded{subcaption}{}{\usepackage{subcaption}}
\makeatother
\makeatletter
\@ifpackageloaded{fontawesome5}{}{\usepackage{fontawesome5}}
\makeatother

% From https://tex.stackexchange.com/a/645996/211326
%%% apa7 doesn't want to add appendix section titles in the toc
%%% let's make it do it
\makeatletter
\xpatchcmd{\appendix}
  {\par}
  {\addcontentsline{toc}{section}{\@currentlabelname}\par}
  {}{}
\makeatother

%% Disable longtable counter
%% https://tex.stackexchange.com/a/248395/211326

\usepackage{etoolbox}

\makeatletter
\patchcmd{\LT@caption}
  {\bgroup}
  {\bgroup\global\LTpatch@captiontrue}
  {}{}
\patchcmd{\longtable}
  {\par}
  {\par\global\LTpatch@captionfalse}
  {}{}
\apptocmd{\endlongtable}
  {\ifLTpatch@caption\else\addtocounter{table}{-1}\fi}
  {}{}
\newif\ifLTpatch@caption
\makeatother

\begin{document}

\maketitle


\hypertarget{toc}{}
\tableofcontents
\newpage
\section[Introduction]{Introducción a la programación con Python}

\setcounter{secnumdepth}{3}

\setlength\LTleft{0pt}




\section{Introducción a la programación con
Python}\label{introducciuxf3n-a-la-programaciuxf3n-con-python}

\begin{quote}
Esta guía es una introducción \textbf{pura} al lenguaje Python,
construida exclusivamente sobre la documentación oficial: el
\emph{Python Tutorial} (docs.python.org/3/tutorial) y la página oficial
de descargas (python.org/downloads). Trata \textbf{únicamente} qué es
Python, su filosofía de diseño, cómo instalarlo y cómo invocar el
intérprete (modo interactivo y modo script). \textbf{Deliberadamente no
cubre} variables, expresiones, statements, objetos, ejecución
condicional, iteraciones, funciones, dataframes ni machine learning ---
esos temas se abordan en guías posteriores y separadas, siguiendo tu
propia secuencia ya establecida (02 a 11). Esta guía es, en ese sentido,
el ``capítulo 0'': el terreno que hay que pisar antes de escribir la
primera línea de código con sentido.
\end{quote}

\begin{center}\rule{0.5\linewidth}{0.5pt}\end{center}

\subsection{Tabla de contenidos}\label{tabla-de-contenidos}

\begin{enumerate}
\def\labelenumi{\arabic{enumi}.}
\tightlist
\item
  \hyperref[1-quuxe9-es-python-definiciuxf3n-oficial]{¿Qué es Python?
  Definición oficial}
\item
  \hyperref[2-por-quuxe9-existe-python-la-motivaciuxf3n-original]{¿Por
  qué existe Python? La motivación original}
\item
  \hyperref[3-caracteruxedsticas-que-distinguen-a-python-como-lenguaje]{Características
  que distinguen a Python como lenguaje}
\item
  \hyperref[4-por-quuxe9-los-programas-en-python-son-compactos-y-legibles]{Por
  qué los programas en Python son compactos y legibles}
\item
  \hyperref[5-extensibilidad-python-como-lenguaje-de-extensiuxf3n]{Extensibilidad:
  Python como lenguaje de extensión}
\item
  \hyperref[6-origen-del-nombre-python]{Origen del nombre ``Python''}
\item
  \hyperref[7-estructura-oficial-del-tutorial-de-python]{Estructura
  oficial del tutorial de Python}
\item
  \hyperref[8-versiones-de-python-y-ciclo-de-vida-oficial]{Versiones de
  Python y ciclo de vida oficial}
\item
  \hyperref[9-instalaciuxf3n-de-python-documentaciuxf3n-oficial]{Instalación
  de Python: documentación oficial}
\item
  \hyperref[10-verificaciuxf3n-de-la-integridad-del-instalador]{Verificación
  de la integridad del instalador}
\item
  \hyperref[11-invocar-el-intuxe9rprete-conceptos-generales]{Invocar el
  intérprete: conceptos generales}
\item
  \hyperref[12-el-modo-interactivo]{El modo interactivo}
\item
  \hyperref[13-ediciuxf3n-de-luxednea-de-comandos-en-el-intuxe9rprete]{Edición
  de línea de comandos en el intérprete}
\item
  \hyperref[14-ejecutar-python-como-script-modo-no-interactivo]{Ejecutar
  Python como script (modo no interactivo)}
\item
  \hyperref[15-la-opciuxf3n--c-ejecutar-un-comando-directo]{La opción
  \texttt{-c}: ejecutar un comando directo}
\item
  \hyperref[16-la-opciuxf3n--m-ejecutar-un-muxf3dulo-como-script]{La
  opción \texttt{-m}: ejecutar un módulo como script}
\item
  \hyperref[17-la-opciuxf3n--i-modo-interactivo-despuuxe9s-de-un-script]{La
  opción \texttt{-i}: modo interactivo después de un script}
\item
  \hyperref[18-paso-de-argumentos-sysargv]{Paso de argumentos:
  \texttt{sys.argv}}
\item
  \hyperref[19-el-intuxe9rprete-y-su-entorno-codificaciuxf3n-del-cuxf3digo-fuente]{El
  intérprete y su entorno: codificación del código fuente}
\item
  \hyperref[20-la-luxednea-shebang-en-unixlinux]{La línea shebang en
  Unix/Linux}
\item
  \hyperref[21-salir-del-intuxe9rprete]{Salir del intérprete}
\item
  \hyperref[22-comentarios-en-python]{Comentarios en Python}
\item
  \hyperref[23-el-intuxe9rprete-como-calculadora-una-vista-previa-controlada]{El
  intérprete como calculadora: una vista previa controlada}
\item
  \hyperref[24-pythonstartup-personalizar-el-inicio-del-modo-interactivo]{PYTHONSTARTUP:
  personalizar el inicio del modo interactivo}
\item
  \hyperref[25-glosario-muxednimo-para-esta-etapa]{Glosario mínimo para
  esta etapa}
\item
  \hyperref[26-quuxe9-sigue-hoja-de-ruta-de-las-pruxf3ximas-guuxedas]{Qué
  sigue: hoja de ruta de las próximas guías}
\item
  \hyperref[27-referencias-oficiales]{Referencias oficiales}
\end{enumerate}

\begin{center}\rule{0.5\linewidth}{0.5pt}\end{center}

\subsection{1. ¿Qué es Python? Definición
oficial}\label{quuxe9-es-python-definiciuxf3n-oficial}

Antes de instalar nada, conviene fijar con precisión qué es Python según
su propia documentación, en vez de partir de descripciones de terceros.

El \emph{Python Tutorial} oficial lo resume así: Python es, ante todo,
un lenguaje pensado para \textbf{automatizar tareas} que surgen del
trabajo cotidiano con computadoras --- por ejemplo, hacer búsqueda y
reemplazo sobre un gran número de archivos de texto, o renombrar y
reorganizar un conjunto de archivos de fotos de forma compleja. También
sirve para escribir una base de datos pequeña y personalizada, una
aplicación gráfica (GUI) especializada, o un juego sencillo.

La documentación es explícita sobre dónde se ubica Python frente a otras
alternativas: comparado con scripts de shell de Unix o archivos batch de
Windows (buenos para mover archivos y cambiar datos de texto, pero no
aptos para aplicaciones GUI o juegos) y comparado con programas en
C/C++/Java (que pueden tardar mucho tiempo de desarrollo incluso para un
primer borrador), Python ocupa un punto intermedio: \textbf{es simple de
usar pero es un lenguaje de programación real}, con mucha más estructura
y soporte para programas grandes que lo que ofrecen los scripts de shell
o los archivos batch --- y al mismo tiempo, ofrece bastante más
verificación de errores que C.

\begin{center}\rule{0.5\linewidth}{0.5pt}\end{center}

\subsection{2. ¿Por qué existe Python? La motivación
original}\label{por-quuxe9-existe-python-la-motivaciuxf3n-original}

Python nació para resolver una frustración concreta y muy específica del
desarrollo de software profesional. Si trabajas con varias bibliotecas
en C/C++/Java, el ciclo habitual de escribir/compilar/probar/recompilar
puede resultar demasiado lento. Si estás escribiendo un conjunto de
pruebas (\emph{test suite}) para una de esas bibliotecas, escribir el
código de prueba puede volverse una tarea tediosa. O quizás escribiste
un programa que podría beneficiarse de un lenguaje de extensión, pero no
quieres diseñar e implementar un lenguaje completamente nuevo solo para
esa aplicación.

Python, según la documentación oficial, \textbf{es justamente el
lenguaje para esos casos}.

\begin{center}\rule{0.5\linewidth}{0.5pt}\end{center}

\subsection{3. Características que distinguen a Python como
lenguaje}\label{caracteruxedsticas-que-distinguen-a-python-como-lenguaje}

La documentación oficial atribuye a Python las siguientes
características distintivas, descritas con precisión técnica (sin entrar
todavía en cómo se usan, eso corresponde a guías futuras):

\textbf{Lenguaje de muy alto nivel (\emph{very-high-level language})}:
Python tiene tipos de datos de alto nivel incorporados de fábrica, como
arreglos flexibles y diccionarios. Por sus tipos de datos más generales,
Python es aplicable a un dominio de problemas mucho más amplio que Awk o
incluso Perl, aunque muchas cosas son al menos igual de sencillas en
Python que en esos lenguajes.

\textbf{Modularidad}: Python permite dividir un programa en módulos que
pueden reutilizarse en otros programas Python. Viene con una amplia
colección de módulos estándar que se pueden usar como base para tus
propios programas, o como ejemplos para empezar a aprender a programar.
Algunos de esos módulos proveen funcionalidades como E/S de archivos,
llamadas al sistema, sockets, e incluso interfaces a kits de
herramientas de interfaz gráfica de usuario como Tk.

\textbf{Lenguaje interpretado}: Python es un lenguaje interpretado, lo
cual puede ahorrar tiempo considerable durante el desarrollo de un
programa porque no es necesaria ninguna compilación ni enlazado
(\emph{linking}). El intérprete puede usarse de forma interactiva, lo
que facilita experimentar con características del lenguaje, escribir
programas desechables, o probar funciones durante un desarrollo de
programa de abajo hacia arriba (\emph{bottom-up}). También funciona,
según la propia documentación oficial, como una práctica calculadora de
escritorio.

Cada una de estas características ---tipos de datos de alto nivel,
módulos, naturaleza interpretada--- se desarrollará con mucho más
detalle en las guías siguientes (variables y expresiones, objetos,
funciones). Aquí solo se presentan como rasgos definitorios del
lenguaje, tal como los describe su documentación oficial.

\begin{center}\rule{0.5\linewidth}{0.5pt}\end{center}

\subsection{4. Por qué los programas en Python son compactos y
legibles}\label{por-quuxe9-los-programas-en-python-son-compactos-y-legibles}

La documentación oficial explica que Python permite escribir programas
de forma compacta y legible, y que los programas escritos en Python son
típicamente mucho más cortos que programas equivalentes en C, C++ o
Java, por varias razones explícitas:

\begin{itemize}
\tightlist
\item
  Los tipos de datos de alto nivel permiten expresar operaciones
  complejas en una sola sentencia (\emph{statement}).
\item
  El agrupamiento de sentencias se hace mediante \textbf{indentación},
  en vez de llaves de apertura y cierre.
\item
  No son necesarias declaraciones de variables ni de argumentos.
\end{itemize}

Estos tres puntos son, según la propia documentación, la explicación
oficial de por qué el código Python suele ``leerse'' de forma más
directa que el de otros lenguajes de propósito similar. El detalle de
cómo funciona cada uno de ellos (qué es exactamente una sentencia, cómo
opera la indentación en la práctica, cómo se declaran variables sin
declaración explícita) corresponde a la guía 02 (``Variables,
expresiones y statements''), no a esta introducción.

\begin{center}\rule{0.5\linewidth}{0.5pt}\end{center}

\subsection{5. Extensibilidad: Python como lenguaje de
extensión}\label{extensibilidad-python-como-lenguaje-de-extensiuxf3n}

Otro rasgo que la documentación oficial destaca explícitamente es que
Python es \textbf{extensible} (\emph{extensible}): si sabes programar en
C, es sencillo añadir una nueva función o módulo incorporado
(\emph{built-in}) al intérprete, ya sea para realizar operaciones
críticas a la máxima velocidad, o para enlazar programas Python con
bibliotecas que solo estén disponibles en forma binaria (como una
biblioteca gráfica específica de un proveedor). Una vez que te enganchas
de verdad con el lenguaje, puedes enlazar el intérprete de Python dentro
de una aplicación escrita en C y usarlo como lenguaje de extensión o de
comandos para esa aplicación.

Esta característica es la que explica, por ejemplo, por qué bibliotecas
científicas de alto rendimiento (NumPy, pandas, scikit-learn, entre
muchas otras que aparecerán en guías posteriores de este recorrido)
pueden combinar la facilidad de uso de Python con un núcleo de cómputo
escrito en C o C++.

\begin{center}\rule{0.5\linewidth}{0.5pt}\end{center}

\subsection{6. Origen del nombre
``Python''}\label{origen-del-nombre-python}

Un dato que la propia documentación oficial menciona explícitamente, y
anima a repetir: el lenguaje recibe su nombre del programa de la BBC
\emph{Monty Python's Flying Circus}, y \textbf{no tiene nada que ver con
reptiles}. Hacer referencias a sketches de Monty Python en la
documentación no solo está permitido, sino que la documentación oficial
lo anima activamente.

\begin{center}\rule{0.5\linewidth}{0.5pt}\end{center}

\subsection{7. Estructura oficial del tutorial de
Python}\label{estructura-oficial-del-tutorial-de-python}

Antes de entrar en instalación y manejo del intérprete, vale la pena
tener clara la hoja de ruta completa que sigue el tutorial oficial de
Python, tal como la describe su propia introducción: el tutorial
presenta varias funcionalidades del lenguaje y del sistema a través de
ejemplos, comenzando por expresiones simples, sentencias y tipos de
datos, pasando luego por funciones y módulos, y finalmente abordando
conceptos más avanzados como excepciones y clases definidas por el
usuario.

Esta hoja de ruta oficial \textbf{coincide, en su orden general, con la
secuencia que ya seguías} en tus guías 02-11
(variables/expresiones/statements → objetos → condicionales →
iteraciones → funciones → dataframes → predicción/ML). Esta guía de
introducción se ubica, por tanto, exactamente donde la documentación
oficial la ubicaría: antes de todo eso.

La documentación también es honesta sobre el método de aprendizaje
recomendado: dado que la mejor forma de aprender un lenguaje es
usándolo, el tutorial invita a jugar con el intérprete de Python
mientras se lee. El siguiente capítulo del tutorial oficial explica la
mecánica de uso del intérprete --- información que, según la propia
documentación, ``es bastante mundana, pero esencial para poder probar
los ejemplos que se muestran más adelante''. Las secciones 9 en adelante
de esta guía cubren exactamente esa mecánica.

\begin{center}\rule{0.5\linewidth}{0.5pt}\end{center}

\subsection{8. Versiones de Python y ciclo de vida
oficial}\label{versiones-de-python-y-ciclo-de-vida-oficial}

Python sigue un calendario de mantenimiento formal definido en PEPs
(\emph{Python Enhancement Proposals}) específicos para cada serie de
versiones (por ejemplo, PEP 619 para la serie 3.10, PEP 664 para la
3.11). Cada serie pasa por etapas: desarrollo activo, mantenimiento con
correcciones de errores regulares, y finalmente una etapa de ``solo
correcciones de seguridad'' antes de su fin de soporte (\emph{end of
support}).

Al momento de escribir esta guía, según la página oficial de descargas
(python.org/downloads), la serie estable más reciente es \textbf{Python
3.14}, con la serie 3.13 también vigente y recibiendo actualizaciones
regulares; las series 3.9, 3.10 y 3.11 se encuentran ya en la etapa de
``solo correcciones de seguridad'', lo que significa que ya no reciben
correcciones de errores regulares y, en el caso de algunas, ya no se
proveen instaladores binarios nuevos para ellas.

\begin{quote}
\textbf{Recomendación práctica derivada de esto}: para un entorno nuevo
de aprendizaje (como el que construirás en las guías siguientes con
conda/mamba), conviene partir de una versión dentro del rango de
mantenimiento activo (actualmente 3.12, 3.13 o 3.14), evitando series
que ya estén solo en modo de parches de seguridad.
\end{quote}

\begin{center}\rule{0.5\linewidth}{0.5pt}\end{center}

\subsection{9. Instalación de Python: documentación
oficial}\label{instalaciuxf3n-de-python-documentaciuxf3n-oficial}

Aunque ya configuraste Miniconda/Anaconda en guías previas ---que
incluyen su propia distribución de Python integrada---, esta sección
documenta la instalación de Python \textbf{directamente desde la fuente
oficial del lenguaje} (python.org), independiente de cualquier
distribución de terceros, siguiendo la página oficial de descargas.

\subsubsection{9.1. Dónde descargar}\label{duxf3nde-descargar}

La fuente oficial y única recomendada por el propio proyecto Python es
\textbf{python.org/downloads}. Esa página lista, para cada serie de
versión, su estado de mantenimiento, fecha de primer lanzamiento, fecha
de fin de soporte, y el PEP que rige su calendario de lanzamientos.

\subsubsection{9.2. Windows}\label{windows}

python.org/downloads/windows ofrece actualmente el \textbf{Python
install manager} como mecanismo recomendado de instalación y gestión de
versiones de Python en Windows, además de los instaladores tradicionales
\texttt{.exe} para cada versión específica (por ejemplo, Python 3.14.x,
Python 3.13.x). Cada entrada de versión indica explícitamente
compatibilidad de sistema operativo --- por ejemplo, la documentación
oficial advierte que ciertas versiones recientes de Python 3.13
\textbf{no pueden usarse en Windows 7 o anterior}.

\subsubsection{9.3. Verificación de firma digital en Windows y
macOS}\label{verificaciuxf3n-de-firma-digital-en-windows-y-macos}

La documentación oficial detalla un mecanismo de verificación de
autenticidad específico por plataforma:

\begin{itemize}
\tightlist
\item
  \textbf{Windows}: los ejecutables de instalación están firmados
  digitalmente. Esto puede verificarse viendo las propiedades del
  archivo ejecutable, en la pestaña de Firmas Digitales (\emph{Digital
  Signatures}), y confirmando el nombre del firmante. El sujeto completo
  del certificado oficial es
  \texttt{CN\ =\ Python\ Software\ Foundation,\ O\ =\ Python\ Software\ Foundation,\ L\ =\ Beaverton,\ S\ =\ Oregon,\ C\ =\ US},
  y desde el 14 de octubre de 2024 la autoridad certificadora es
  Microsoft Identity Verification Root Certificate Authority
  (anteriormente, los certificados eran emitidos por DigiCert). La
  documentación aclara que algunos ejecutables pueden no estar firmados
  ---notablemente, el comando \texttt{pip} por defecto---, ya que no se
  construyen como parte de Python sino que se incluyen desde bibliotecas
  de terceros.
\item
  \textbf{macOS}: los paquetes instaladores descargables desde
  python.org están firmados con un certificado Apple Developer ID
  Installer. Desde Python 3.11.4 y 3.12.0b1 (23 de mayo de 2023), los
  paquetes instaladores de release se firman con certificados emitidos a
  la Python Software Foundation.
\end{itemize}

\subsubsection{9.4. Linux: nota sobre el alcance de esta
guía}\label{linux-nota-sobre-el-alcance-de-esta-guuxeda}

En Kubuntu y Arch Linux, una versión de Python suele venir preinstalada
como parte del sistema, o se gestiona vía el gestor de paquetes de la
distribución (\texttt{apt}/\texttt{pacman}) o, como ya documentaste en
guías previas, vía Miniconda/Anaconda. La instalación de Python ``puro''
desde fuente en Linux (compilación desde el código fuente publicado en
python.org) es un procedimiento avanzado que excede el alcance de esta
introducción; la documentación oficial recomienda, para la mayoría de
usuarios de Linux, usar el gestor de paquetes de la distribución o una
distribución como Anaconda/Miniconda --- que es exactamente la vía que
ya tienes configurada.

\begin{center}\rule{0.5\linewidth}{0.5pt}\end{center}

\subsection{10. Verificación de la integridad del
instalador}\label{verificaciuxf3n-de-la-integridad-del-instalador}

Más allá de la firma digital descrita en la sección anterior, la
práctica general (consistente con lo ya documentado en tu guía de
instalación de Miniconda/Anaconda) es verificar que el archivo
descargado no haya sido alterado ni corrompido, comparando su hash con
el valor publicado oficialmente junto al instalador en la página de
descargas correspondiente, antes de ejecutarlo.

\begin{center}\rule{0.5\linewidth}{0.5pt}\end{center}

\subsection{11. Invocar el intérprete: conceptos
generales}\label{invocar-el-intuxe9rprete-conceptos-generales}

A partir de aquí, la guía sigue de cerca el capítulo oficial ``Using the
Python Interpreter'' (\emph{Usando el intérprete de Python}) del
tutorial.

\subsubsection{11.1. Ubicación habitual del
ejecutable}\label{ubicaciuxf3n-habitual-del-ejecutable}

Según la documentación oficial, el intérprete de Python suele instalarse
como \texttt{/usr/local/bin/python3.14} (ajustando el número de versión
según corresponda) en las máquinas donde está disponible; agregar
\texttt{/usr/local/bin} a la ruta de búsqueda de tu shell de Unix
permite iniciarlo escribiendo el comando:

\begin{Shaded}
\begin{Highlighting}[]
\ExtensionTok{python3.14}
\end{Highlighting}
\end{Shaded}

La documentación aclara que, dado que la elección del directorio donde
vive el intérprete es una opción de instalación, son posibles otras
ubicaciones (por ejemplo, \texttt{/usr/local/python} es una alternativa
popular); conviene verificar con la administración del sistema o la
documentación de la propia distribución si hay dudas.

\begin{quote}
\textbf{Nota oficial relevante para Kubuntu/Arch}: en Unix, el
intérprete de Python 3.x \textbf{no se instala por defecto} con el
ejecutable llamado simplemente \texttt{python}, precisamente para no
entrar en conflicto con una instalación simultánea de un ejecutable de
Python 2.x. Esta es la razón oficial por la que, en muchos sistemas
Linux, el comando correcto suele ser \texttt{python3} y no
\texttt{python}.
\end{quote}

\subsubsection{11.2. En Windows}\label{en-windows}

En máquinas Windows donde instalaste Python desde la Microsoft Store, el
comando \texttt{python3.14} (o el número de versión correspondiente)
estará disponible. Si tienes instalado el lanzador \texttt{py.exe},
puedes usar el comando \texttt{py}. La documentación oficial remite a la
sección sobre el \emph{Python install manager} para conocer otras formas
de iniciar Python en Windows.

\subsubsection{11.3. Comportamiento general, similar al shell de
Unix}\label{comportamiento-general-similar-al-shell-de-unix}

El intérprete opera de forma parecida al shell de Unix: cuando se le
llama con la entrada estándar conectada a un dispositivo tty, lee y
ejecuta comandos de forma interactiva; cuando se le llama con un
argumento de nombre de archivo, o con un archivo como entrada estándar,
lee y ejecuta un \emph{script} desde ese archivo.

\begin{center}\rule{0.5\linewidth}{0.5pt}\end{center}

\subsection{12. El modo interactivo}\label{el-modo-interactivo}

Cuando los comandos se leen desde un tty, se dice que el intérprete está
en \textbf{modo interactivo} (\emph{interactive mode}). En este modo,
solicita el siguiente comando con el \textbf{prompt primario},
normalmente tres signos de mayor que
(\texttt{\textgreater{}\textgreater{}\textgreater{}}); para líneas de
continuación, solicita con el \textbf{prompt secundario}, por defecto
tres puntos (\texttt{...}).

El intérprete imprime un mensaje de bienvenida indicando su número de
versión y un aviso de copyright antes de mostrar el primer prompt:

\begin{verbatim}
$ python3.14
Python 3.14 (default, April 4 2024, 09:25:04)
[GCC 10.2.0] on linux
Type "help", "copyright", "credits" or "license" for more information.
>>>
\end{verbatim}

Las líneas de continuación son necesarias al introducir una construcción
de varias líneas. El ejemplo que la propia documentación oficial usa
para ilustrarlo (sin que esto sea todavía una explicación de
condicionales, que corresponde a la guía 04) es el siguiente:

\begin{Shaded}
\begin{Highlighting}[]
\OperatorTok{\textgreater{}\textgreater{}\textgreater{}}\NormalTok{ the\_world\_is\_flat }\OperatorTok{=} \VariableTok{True}
\OperatorTok{\textgreater{}\textgreater{}\textgreater{}} \ControlFlowTok{if}\NormalTok{ the\_world\_is\_flat:}
\NormalTok{...     }\BuiltInTok{print}\NormalTok{(}\StringTok{"Be careful not to fall off!"}\NormalTok{)}
\NormalTok{...}
\NormalTok{Be careful }\KeywordTok{not}\NormalTok{ to fall off}\OperatorTok{!}
\end{Highlighting}
\end{Shaded}

Para esta guía, lo relevante de ese ejemplo no es \emph{qué hace} el
\texttt{if} (eso es tema de una guía futura), sino \emph{cómo se ve} la
interacción con el intérprete: el prompt primario
\texttt{\textgreater{}\textgreater{}\textgreater{}} para la primera
línea, el prompt secundario \texttt{...} para las líneas de
continuación, y una línea en blanco con solo \texttt{...} para indicarle
al intérprete que el bloque de varias líneas ha terminado.

\begin{center}\rule{0.5\linewidth}{0.5pt}\end{center}

\subsection{13. Edición de línea de comandos en el
intérprete}\label{ediciuxf3n-de-luxednea-de-comandos-en-el-intuxe9rprete}

Las funciones de edición de línea del intérprete incluyen edición
interactiva, sustitución de historial, y autocompletado de código en la
mayoría de los sistemas, según la documentación oficial.

La forma más rápida de comprobar si tu sistema soporta edición de línea
de comandos es escribir una palabra en el prompt de Python y luego
presionar la flecha izquierda (o \texttt{Control-b}). Si el cursor se
mueve, tienes edición de línea de comandos disponible. Si no ocurre nada
visible, o aparece una secuencia como \texttt{\^{}{[}{[}D} o
\texttt{\^{}B}, la edición de línea de comandos no está disponible, y
solo podrás usar la tecla de retroceso (\emph{backspace}) para borrar
caracteres de la línea actual.

\begin{center}\rule{0.5\linewidth}{0.5pt}\end{center}

\subsection{14. Ejecutar Python como script (modo no
interactivo)}\label{ejecutar-python-como-script-modo-no-interactivo}

Cuando se le da al intérprete un argumento con el nombre de un archivo,
lee y ejecuta ese archivo como un \emph{script}, en vez de entrar en
modo interactivo. Esta es la forma habitual de ejecutar programas Python
guardados en un archivo \texttt{.py}:

\begin{Shaded}
\begin{Highlighting}[]
\ExtensionTok{python3}\NormalTok{ mi\_programa.py}
\end{Highlighting}
\end{Shaded}

Esta guía no desarrolla todavía el contenido que iría dentro de
\texttt{mi\_programa.py} ---eso corresponde a partir de la guía 02 en
adelante---, pero sí establece la mecánica de invocación, que es
información de esta etapa ``previa al código'' según la propia
estructura del tutorial oficial.

\begin{center}\rule{0.5\linewidth}{0.5pt}\end{center}

\subsection{\texorpdfstring{15. La opción \texttt{-c}: ejecutar un
comando
directo}{15. La opción -c: ejecutar un comando directo}}\label{la-opciuxf3n--c-ejecutar-un-comando-directo}

Una segunda forma de iniciar el intérprete es
\texttt{python\ -c\ comando\ {[}arg{]}\ ...}, que ejecuta la(s)
sentencia(s) en \emph{comando}, de forma análoga a la opción \texttt{-c}
del shell. Dado que las sentencias de Python suelen contener espacios u
otros caracteres especiales para el shell, la documentación oficial
recomienda encerrar \emph{comando} entre comillas en su totalidad.

\begin{Shaded}
\begin{Highlighting}[]
\ExtensionTok{python3} \AttributeTok{{-}c} \StringTok{"print(\textquotesingle{}hola\textquotesingle{})"}
\end{Highlighting}
\end{Shaded}

\begin{center}\rule{0.5\linewidth}{0.5pt}\end{center}

\subsection{\texorpdfstring{16. La opción \texttt{-m}: ejecutar un
módulo como
script}{16. La opción -m: ejecutar un módulo como script}}\label{la-opciuxf3n--m-ejecutar-un-muxf3dulo-como-script}

Algunos módulos de Python también son útiles como scripts. Pueden
invocarse usando \texttt{python\ -m\ módulo\ {[}arg{]}\ ...}, lo cual
ejecuta el archivo fuente correspondiente a \emph{módulo} como si
hubieras escrito su nombre completo en la línea de comandos.

\begin{Shaded}
\begin{Highlighting}[]
\ExtensionTok{python3} \AttributeTok{{-}m}\NormalTok{ venv mi\_entorno}
\end{Highlighting}
\end{Shaded}

(el ejemplo anterior es solo ilustrativo de la sintaxis \texttt{-m}; el
manejo de entornos virtuales con \texttt{venv} no es tema de esta guía,
ya que tu flujo de trabajo usa conda/mamba, documentado en guías
previas).

\begin{center}\rule{0.5\linewidth}{0.5pt}\end{center}

\subsection{\texorpdfstring{17. La opción \texttt{-i}: modo interactivo
después de un
script}{17. La opción -i: modo interactivo después de un script}}\label{la-opciuxf3n--i-modo-interactivo-despuuxe9s-de-un-script}

Cuando se usa un archivo de script, a veces es útil poder ejecutar el
script y luego entrar en modo interactivo. Esto puede hacerse pasando
\texttt{-i} antes del script:

\begin{Shaded}
\begin{Highlighting}[]
\ExtensionTok{python3} \AttributeTok{{-}i}\NormalTok{ mi\_programa.py}
\end{Highlighting}
\end{Shaded}

Todas las opciones de línea de comandos están descritas, según remite la
documentación oficial, en la sección ``Command line and environment'' de
la referencia de Python.

\begin{center}\rule{0.5\linewidth}{0.5pt}\end{center}

\subsection{\texorpdfstring{18. Paso de argumentos:
\texttt{sys.argv}}{18. Paso de argumentos: sys.argv}}\label{paso-de-argumentos-sys.argv}

Cuando el intérprete los reconoce, el nombre del script y los argumentos
adicionales que le siguen se convierten en una lista de cadenas de texto
y se asignan a la variable \texttt{argv} dentro del módulo \texttt{sys}.
Se puede acceder a esa lista ejecutando \texttt{import\ sys}.

Algunos detalles precisos de la documentación oficial sobre esta
variable:

\begin{itemize}
\tightlist
\item
  La longitud de la lista es de al menos uno; cuando no se da ningún
  script ni argumentos, \texttt{sys.argv{[}0{]}} es una cadena vacía.
\item
  Cuando el nombre del script se da como
  \texttt{\textquotesingle{}-\textquotesingle{}} (es decir, entrada
  estándar), \texttt{sys.argv{[}0{]}} se fija como
  \texttt{\textquotesingle{}-\textquotesingle{}}.
\item
  Cuando se usa \texttt{-c\ comando}, \texttt{sys.argv{[}0{]}} se fija
  como \texttt{\textquotesingle{}-c\textquotesingle{}}.
\item
  Cuando se usa \texttt{-m\ módulo}, \texttt{sys.argv{[}0{]}} se fija
  con el nombre completo del módulo localizado.
\item
  Las opciones que aparecen después de \texttt{-c\ comando} o
  \texttt{-m\ módulo} no son consumidas por el procesamiento de opciones
  del propio intérprete de Python, sino que se dejan en
  \texttt{sys.argv} para que las maneje el comando o módulo
  correspondiente.
\end{itemize}

\begin{quote}
Esta sección se documenta aquí por completitud y rigor (es parte literal
del capítulo oficial sobre el intérprete), aunque su uso práctico pleno
requiere conceptos de listas y módulos que se desarrollarán en guías
posteriores.
\end{quote}

\begin{center}\rule{0.5\linewidth}{0.5pt}\end{center}

\subsection{19. El intérprete y su entorno: codificación del código
fuente}\label{el-intuxe9rprete-y-su-entorno-codificaciuxf3n-del-cuxf3digo-fuente}

Por defecto, los archivos fuente de Python se tratan como codificados en
\textbf{UTF-8}. En esa codificación, pueden usarse simultáneamente
caracteres de la mayoría de los idiomas del mundo en literales de
cadena, identificadores y comentarios --- aunque la biblioteca estándar
solo usa caracteres ASCII para identificadores, una convención que
cualquier código portable debería seguir. Para mostrar correctamente
todos esos caracteres, tu editor debe reconocer que el archivo está en
UTF-8 y debe usar una fuente que soporte todos los caracteres del
archivo.

\begin{quote}
\textbf{Nota práctica para tu entorno}: dado que trabajas con Neovim
sobre Kubuntu/Arch, UTF-8 es generalmente la codificación por defecto en
ambos sistemas, por lo que normalmente no necesitarás declarar nada
adicional para escribir comentarios o cadenas en español con tildes y
eñes.
\end{quote}

\subsubsection{19.1. Declarar una codificación distinta a la
predeterminada}\label{declarar-una-codificaciuxf3n-distinta-a-la-predeterminada}

Para declarar una codificación distinta de la predeterminada, debe
agregarse una línea de comentario especial como \textbf{primera línea}
del archivo, con la siguiente sintaxis:

\begin{Shaded}
\begin{Highlighting}[]
\CommentTok{\# {-}*{-} coding: encoding {-}*{-}}
\end{Highlighting}
\end{Shaded}

donde \emph{encoding} es uno de los códecs (\texttt{codecs}) válidos
soportados por Python. Por ejemplo, para declarar que se use la
codificación Windows-1252, la primera línea del archivo fuente debería
ser:

\begin{Shaded}
\begin{Highlighting}[]
\CommentTok{\# {-}*{-} coding: cp1252 {-}*{-}}
\end{Highlighting}
\end{Shaded}

\begin{center}\rule{0.5\linewidth}{0.5pt}\end{center}

\subsection{20. La línea shebang en
Unix/Linux}\label{la-luxednea-shebang-en-unixlinux}

Una excepción a la regla de ``primera línea'' es cuando el código fuente
comienza con una \textbf{línea shebang} de Unix. En ese caso, la
declaración de codificación debe agregarse como \textbf{segunda línea}
del archivo. Ejemplo oficial:

\begin{Shaded}
\begin{Highlighting}[]
\CommentTok{\#!/usr/bin/env python3}
\CommentTok{\# {-}*{-} coding: cp1252 {-}*{-}}
\end{Highlighting}
\end{Shaded}

La línea shebang (\texttt{\#!/usr/bin/env\ python3}) es lo que permite,
en Kubuntu y Arch Linux, hacer un script ejecutable directamente
(\texttt{./mi\_script.py}) sin tener que invocar \texttt{python3}
explícitamente, siempre que el archivo tenga el permiso de ejecución
correspondiente (otorgado, por ejemplo, con \texttt{chmod\ +x}).

\begin{center}\rule{0.5\linewidth}{0.5pt}\end{center}

\subsection{21. Salir del intérprete}\label{salir-del-intuxe9rprete}

Escribir un carácter de fin de archivo (\emph{end-of-file}) ---
\texttt{Control-D} en Unix, \texttt{Control-Z} en Windows --- en el
prompt primario hace que el intérprete salga con un estado de salida
cero. Si eso no funciona, puedes salir del intérprete escribiendo el
siguiente comando:

\begin{Shaded}
\begin{Highlighting}[]
\NormalTok{quit()}
\end{Highlighting}
\end{Shaded}

\begin{center}\rule{0.5\linewidth}{0.5pt}\end{center}

\subsection{22. Comentarios en Python}\label{comentarios-en-python}

Los comentarios son una de las pocas piezas de sintaxis que sí
corresponde introducir aquí, porque son necesarios para leer y escribir
cualquier ejemplo de código desde el primer momento, incluso antes de
cubrir variables o expresiones en profundidad.

Según la documentación oficial: \textbf{los comentarios en Python
comienzan con el carácter numeral (\texttt{\#}) y se extienden hasta el
final de la línea física}. Un comentario puede aparecer al inicio de una
línea, o después de espacio en blanco o código, pero \textbf{no dentro
de un literal de cadena de texto}. Un carácter \texttt{\#} dentro de un
literal de cadena es simplemente un carácter numeral, no el inicio de un
comentario.

Ejemplo oficial textual:

\begin{Shaded}
\begin{Highlighting}[]
\CommentTok{\# this is the first comment}
\NormalTok{spam }\OperatorTok{=} \DecValTok{1}  \CommentTok{\# and this is the second comment}
          \CommentTok{\# ... and now a third!}
\NormalTok{text }\OperatorTok{=} \StringTok{"\# This is not a comment because it\textquotesingle{}s inside quotes."}
\end{Highlighting}
\end{Shaded}

Dado que los comentarios sirven para clarificar el código y no son
interpretados por Python, la documentación oficial señala que pueden
omitirse al transcribir los ejemplos del propio tutorial --- una
observación útil para cuando, en guías futuras, copies y adaptes
ejemplos de código.

\begin{center}\rule{0.5\linewidth}{0.5pt}\end{center}

\subsection{23. El intérprete como calculadora: una vista previa
controlada}\label{el-intuxe9rprete-como-calculadora-una-vista-previa-controlada}

El tutorial oficial, inmediatamente después de explicar cómo invocar el
intérprete, presenta a Python ``actuando como una simple calculadora'':
puedes escribir una expresión y el intérprete escribirá su valor. Se
incluye aquí únicamente como evidencia de la naturaleza interactiva del
intérprete ---no como introducción a expresiones o tipos numéricos, que
se desarrollarán con rigor en la guía 02---, mostrando el ejemplo
textual oficial:

\begin{Shaded}
\begin{Highlighting}[]
\OperatorTok{\textgreater{}\textgreater{}\textgreater{}} \DecValTok{2} \OperatorTok{+} \DecValTok{2}
\DecValTok{4}
\OperatorTok{\textgreater{}\textgreater{}\textgreater{}} \DecValTok{50} \OperatorTok{{-}} \DecValTok{5}\OperatorTok{*}\DecValTok{6}
\DecValTok{20}
\OperatorTok{\textgreater{}\textgreater{}\textgreater{}}\NormalTok{ (}\DecValTok{50} \OperatorTok{{-}} \DecValTok{5}\OperatorTok{*}\DecValTok{6}\NormalTok{) }\OperatorTok{/} \DecValTok{4}
\FloatTok{5.0}
\OperatorTok{\textgreater{}\textgreater{}\textgreater{}} \DecValTok{8} \OperatorTok{/} \DecValTok{5}  \CommentTok{\# la división siempre devuelve un número de punto flotante}
\FloatTok{1.6}
\end{Highlighting}
\end{Shaded}

Lo único que esta guía retiene de este ejemplo, deliberadamente, es la
mecánica de interacción (escribes después de
\texttt{\textgreater{}\textgreater{}\textgreater{}}, el intérprete
responde en la línea siguiente sin prompt) y el hecho documentado de que
la división con \texttt{/} siempre devuelve un valor de tipo flotante.
El desarrollo completo de tipos de datos, operadores aritméticos,
precedencia, y todo lo demás relacionado a expresiones queda, tal como
pediste, para la guía 02 (``Variables, expresiones y statements con
Python'').

\begin{center}\rule{0.5\linewidth}{0.5pt}\end{center}

\subsection{24. PYTHONSTARTUP: personalizar el inicio del modo
interactivo}\label{pythonstartup-personalizar-el-inicio-del-modo-interactivo}

Cuando usas Python de forma interactiva, suele ser útil tener algunos
comandos estándar ejecutándose automáticamente cada vez que se inicia el
intérprete. Esto puede lograrse fijando una variable de entorno llamada
\texttt{PYTHONSTARTUP} con el nombre de un archivo que contenga esos
comandos de inicio --- de forma similar a la funcionalidad
\texttt{.profile} de los shells de Unix.

Puntos oficiales relevantes sobre este mecanismo:

\begin{itemize}
\tightlist
\item
  Este archivo \textbf{solo se lee en sesiones interactivas}, no cuando
  Python lee comandos desde un script.
\item
  Se ejecuta en el mismo espacio de nombres donde se ejecutan los
  comandos interactivos, de modo que los objetos que define o importa
  pueden usarse sin calificación adicional en la sesión interactiva.
\item
  También es posible cambiar los prompts \texttt{sys.ps1} y
  \texttt{sys.ps2} desde ese archivo.
\end{itemize}

\begin{quote}
Esta funcionalidad se documenta aquí por rigor y completitud frente a la
fuente oficial, aunque su configuración práctica (qué poner dentro de
ese archivo) requiere conceptos ---importación de módulos, asignación de
variables--- que se desarrollarán en guías posteriores.
\end{quote}

\begin{center}\rule{0.5\linewidth}{0.5pt}\end{center}

\subsection{25. Glosario mínimo para esta
etapa}\label{glosario-muxednimo-para-esta-etapa}

Para cerrar esta introducción sin adelantar contenido de guías futuras,
este glosario recoge únicamente los términos que aparecieron en esta
guía, con la definición estrictamente ligada al alcance ya cubierto:

{\def\LTcaptype{none} % do not increment counter
\begin{longtable}[]{@{}
  >{\raggedright\arraybackslash}p{(\linewidth - 2\tabcolsep) * \real{0.5000}}
  >{\raggedright\arraybackslash}p{(\linewidth - 2\tabcolsep) * \real{0.5000}}@{}}
\toprule\noalign{}
\begin{minipage}[b]{\linewidth}\raggedright
Término
\end{minipage} & \begin{minipage}[b]{\linewidth}\raggedright
Definición en el alcance de esta guía
\end{minipage} \\
\midrule\noalign{}
\endhead
\bottomrule\noalign{}
\endlastfoot
Intérprete & Programa que lee y ejecuta código Python, ya sea de forma
interactiva o desde un archivo de script \\
Modo interactivo & Modo del intérprete activado cuando lee comandos
desde un tty; usa los prompts
\texttt{\textgreater{}\textgreater{}\textgreater{}} y \texttt{...} \\
Script & Archivo de texto con extensión \texttt{.py} que el intérprete
lee y ejecuta de forma no interactiva \\
Prompt primario & El símbolo
\texttt{\textgreater{}\textgreater{}\textgreater{}} que indica que el
intérprete espera una nueva instrucción \\
Prompt secundario & El símbolo \texttt{...} que indica que el intérprete
espera la continuación de una instrucción de varias líneas \\
Comentario & Texto que comienza con \texttt{\#} y se extiende hasta el
final de la línea física; ignorado por el intérprete \\
Shebang & Línea inicial de un script Unix
(\texttt{\#!/usr/bin/env\ python3}) que indica qué intérprete debe
ejecutarlo \\
\texttt{sys.argv} & Lista de cadenas de texto con el nombre del script y
los argumentos pasados en la línea de comandos \\
\end{longtable}
}

\begin{center}\rule{0.5\linewidth}{0.5pt}\end{center}

\subsection{26. Qué sigue: hoja de ruta de las próximas
guías}\label{quuxe9-sigue-hoja-de-ruta-de-las-pruxf3ximas-guuxedas}

Siguiendo exactamente la secuencia que ya estableciste, y que coincide
con la estructura del propio tutorial oficial de Python (sección 7 de
esta guía), las siguientes guías cubrirán, en este orden:

\begin{enumerate}
\def\labelenumi{\arabic{enumi}.}
\tightlist
\item
  \textbf{Variables, expresiones y statements con Python} --- tipos de
  datos numéricos y de texto, operadores, precedencia, asignación,
  convenciones de nombres.
\item
  \textbf{Objetos de Python} --- el modelo de objetos del lenguaje.
\item
  \textbf{Ejecución condicional con Python} --- \texttt{if},
  \texttt{elif}, \texttt{else}, operadores de comparación y lógicos.
\item
  \textbf{Iteraciones con Python} --- \texttt{while}, \texttt{for},
  control de bucles.
\item
  \textbf{Funciones con Python} --- definición, parámetros, retorno de
  valores, ámbito.
\item
  \textbf{Dataframes con Python} --- introducción a pandas.
\item
  \textbf{Predicción y métricas de performance con Python}.
\item
  \textbf{Métodos de Machine Learning para clasificación con Python}.
\item
  \textbf{Métodos de Machine Learning para regresión con Python}.
\item
  \textbf{Validación cruzada y composición del modelo con Python}.
\end{enumerate}

Esta guía no se adelanta a ninguno de esos contenidos; su función es
exclusivamente la de establecer, con la documentación oficial como única
fuente, qué es Python y cómo se invoca el intérprete, antes de empezar a
escribir código con sentido semántico.

\begin{center}\rule{0.5\linewidth}{0.5pt}\end{center}

\subsection{27. Referencias oficiales}\label{referencias-oficiales}

\begin{itemize}
\tightlist
\item
  \textbf{Python Tutorial --- Índice general}:
  https://docs.python.org/3/tutorial/index.html
\item
  \textbf{Python Tutorial --- Cap. 1: Whetting Your Appetite}:
  https://docs.python.org/3/tutorial/appetite.html
\item
  \textbf{Python Tutorial --- Cap. 2: Using the Python Interpreter}:
  https://docs.python.org/3/tutorial/interpreter.html
\item
  \textbf{Python Tutorial --- Cap. 3: An Informal Introduction to
  Python} (solo sección de comentarios usada en esta guía):
  https://docs.python.org/3/tutorial/introduction.html
\item
  \textbf{Python --- Página oficial de descargas}:
  https://www.python.org/downloads/
\item
  \textbf{Python --- Descargas para Windows}:
  https://www.python.org/downloads/windows/
\item
  \textbf{PEP 619 --- Python 3.10 Release Schedule}:
  https://peps.python.org/pep-0619/
\item
  \textbf{PEP 664 --- Python 3.11 Release Schedule}:
  https://peps.python.org/pep-0664/
\item
  \textbf{Python --- Glosario oficial}:
  https://docs.python.org/3/glossary.html
\item
  \textbf{Python --- Command line and environment (referencia completa
  de opciones del intérprete)}:
  https://docs.python.org/3/using/cmdline.html
\end{itemize}

\begin{center}\rule{0.5\linewidth}{0.5pt}\end{center}

\textbf{Edison Achalma} Economista --- Universidad Nacional de San
Cristóbal de Huamanga (UNSCH) Ayacucho, Perú ORCID:
\href{https://orcid.org/0000-0001-6996-3364}{0000-0001-6996-3364}

\section{Publicaciones Similares}\label{publicaciones-similares}

Si te interesó este artículo, te recomendamos que explores otros blogs y
recursos relacionados que pueden ampliar tus conocimientos. Aquí te dejo
algunas sugerencias:

\begin{enumerate}
\def\labelenumi{\arabic{enumi}.}
\tightlist
\item
  \href{https://numerus-scriptum.netlify.app/python/2020-06-19-instalacion-de-anaconda/index.pdf}{\faIcon{file-pdf}}
  \href{https://numerus-scriptum.netlify.app/python/2020-06-19-instalacion-de-anaconda}{Instalacion
  De Anaconda}
\item
  \href{https://numerus-scriptum.netlify.app/python/2020-06-20-configurar-entorno-virtual-python-anaconda/index.pdf}{\faIcon{file-pdf}}
  \href{https://numerus-scriptum.netlify.app/python/2020-06-20-configurar-entorno-virtual-python-anaconda}{Configurar
  Entorno Virtual Python Anaconda}
\item
  \href{https://numerus-scriptum.netlify.app/python/2021-04-17-01-introducion-a-la-programacion-con-python/index.pdf}{\faIcon{file-pdf}}
  \href{https://numerus-scriptum.netlify.app/python/2021-04-17-01-introducion-a-la-programacion-con-python}{01
  Introducion A La Programacion Con Python}
\item
  \href{https://numerus-scriptum.netlify.app/python/2021-05-31-02-variables-expresiones-y-statements-con-python/index.pdf}{\faIcon{file-pdf}}
  \href{https://numerus-scriptum.netlify.app/python/2021-05-31-02-variables-expresiones-y-statements-con-python}{02
  Variables Expresiones Y Statements Con Python}
\item
  \href{https://numerus-scriptum.netlify.app/python/2021-06-07-03-objetos-de-python/index.pdf}{\faIcon{file-pdf}}
  \href{https://numerus-scriptum.netlify.app/python/2021-06-07-03-objetos-de-python}{03
  Objetos De Python}
\item
  \href{https://numerus-scriptum.netlify.app/python/2021-06-14-04-ejecucion-condicional-con-python/index.pdf}{\faIcon{file-pdf}}
  \href{https://numerus-scriptum.netlify.app/python/2021-06-14-04-ejecucion-condicional-con-python}{04
  Ejecucion Condicional Con Python}
\item
  \href{https://numerus-scriptum.netlify.app/python/2021-06-21-05-iteraciones-con-python/index.pdf}{\faIcon{file-pdf}}
  \href{https://numerus-scriptum.netlify.app/python/2021-06-21-05-iteraciones-con-python}{05
  Iteraciones Con Python}
\item
  \href{https://numerus-scriptum.netlify.app/python/2021-08-16-06-funciones-con-python/index.pdf}{\faIcon{file-pdf}}
  \href{https://numerus-scriptum.netlify.app/python/2021-08-16-06-funciones-con-python}{06
  Funciones Con Python}
\item
  \href{https://numerus-scriptum.netlify.app/python/2021-08-23-07-dataframes-con-python/index.pdf}{\faIcon{file-pdf}}
  \href{https://numerus-scriptum.netlify.app/python/2021-08-23-07-dataframes-con-python}{07
  Dataframes Con Python}
\item
  \href{https://numerus-scriptum.netlify.app/python/2021-11-29-08-prediccion-y-metrica-de-performance-con-python/index.pdf}{\faIcon{file-pdf}}
  \href{https://numerus-scriptum.netlify.app/python/2021-11-29-08-prediccion-y-metrica-de-performance-con-python}{08
  Prediccion Y Metrica De Performance Con Python}
\item
  \href{https://numerus-scriptum.netlify.app/python/2021-12-06-09-metodos-de-machine-learning-para-clasificacion-con-python/index.pdf}{\faIcon{file-pdf}}
  \href{https://numerus-scriptum.netlify.app/python/2021-12-06-09-metodos-de-machine-learning-para-clasificacion-con-python}{09
  Metodos De Machine Learning Para Clasificacion Con Python}
\item
  \href{https://numerus-scriptum.netlify.app/python/2021-12-13-10-metodos-de-machine-learning-para-regresion-con-python/index.pdf}{\faIcon{file-pdf}}
  \href{https://numerus-scriptum.netlify.app/python/2021-12-13-10-metodos-de-machine-learning-para-regresion-con-python}{10
  Metodos De Machine Learning Para Regresion Con Python}
\item
  \href{https://numerus-scriptum.netlify.app/python/2022-10-31-11-validacion-cruzada-y-composicion-del-modelo-con-python/index.pdf}{\faIcon{file-pdf}}
  \href{https://numerus-scriptum.netlify.app/python/2022-10-31-11-validacion-cruzada-y-composicion-del-modelo-con-python}{11
  Validacion Cruzada Y Composicion Del Modelo Con Python}
\item
  \href{https://numerus-scriptum.netlify.app/python/2025-05-10-visualizacion-de-datos-con-python/index.pdf}{\faIcon{file-pdf}}
  \href{https://numerus-scriptum.netlify.app/python/2025-05-10-visualizacion-de-datos-con-python}{Visualizacion
  De Datos Con Python}
\end{enumerate}

Esperamos que encuentres estas publicaciones igualmente interesantes y
útiles. ¡Disfruta de la lectura!






\end{document}
